\section{Background}
\label{sec:background}

The Bailout Protocol is situated at the intersection of three established research traditions: accountability in distributed AI systems, human-in-the-loop design, and fail-safe engineering in safety-critical domains. It is architecturally instantiated within the \bxthree{} Framework's three-layer cognitive model. This section traces each lineage, identifies the specific gaps that motivate the protocol, and defines the three-layer model against which the protocol operates.

% --------------------------------------------------------------
\subsection{Accountability in Multi-Agent AI Systems}
\label{sec:accountability}

Multi-agent AI systems distribute causal responsibility across networks of autonomous and semi-autonomous nodes in ways that differ fundamentally from both human organizations and traditional software systems. When an agent decomposes a directive, spawns sub-agents, delegates subtasks laterally, or operates asynchronously under partial observability, the chain of responsibility between a human principal and a consequential outcome becomes mediated by multiple machine actors in patterns that are non-linear, opaque, and legally ambiguous~\cite{bansal2019}. This \emph{accountability gap} — the absence of a reliable, traceable pathway from outcome to principal — is among the most significant barriers to deploying AI systems in high-stakes domains.

Classical accountability frameworks borrowed from organizational theory assume a fixed hierarchy of human agents with delineated authority and formal reporting relationships. In software systems, these manifest as role-based access control, audit logs, and approval workflows. These mechanisms were designed for deterministic systems in which all possible states are enumerable at design time. Large language model--based agents introduce states that are not exhaustively predictable, whose behaviors emerge from learned statistical patterns rather than explicit rules, and whose failure modes are neither fully predictable nor fully auditable post hoc~\cite{dijkstra1982}. A node may produce an outcome that is neither clearly authorized nor clearly prohibited within its provisioned scope — a grey zone that binary permission models cannot capture.

The \bxthree{} Framework addresses this through an architectural accountability model in which every node carries an immutable parent pointer and a designated human anchor. The parent pointer defines the chain of delegation; the human anchor ensures that some named principal is architecturally reachable from every node regardless of delegation depth. This is not a governance convention but an architectural constraint: the \bxthree{} state machine does not permit any machine actor to serve as a terminal resolution point for an unresolved exception. The accountability chain is append-only and cryptographically sealed via the Forensic Ledger, enabling complete reconstruction of delegation and exception history for any node at any point in time~\cite{wiener1948}. The relationship is one of necessity: accountability without auditability is intent; auditability without accountability is a record without enforcement. The \bxthree{} Forensic Ledger provides the latter; the Bailout Protocol, operating through the \purpose{} and \bounds{} layers, provides the former.

% --------------------------------------------------------------
\subsection{Human-in-the-Loop Design}
\label{sec:hitl}

Human-in-the-loop (HITL) design positions human agents as integral components of AI-driven workflows rather than as external reviewers who intervene only after a system has reached a conclusion. The canonical motivation is epistemic: AI systems operate within the bounds of their training data and objective functions and cannot reliably reason about novel situations, value trade-offs, or contextual factors that exceed their provisioned epistemic scope~\cite{tristr81}. A secondary motivation — equally pressing in safety-critical deployments — is institutional: many decisions carry legal, financial, or social consequences that cannot be delegated to a machine by institutional design, regardless of empirical performance~\cite{bansal2019}.

Trist's analysis of organizational change established that technological and social systems are co-constituted: introducing a new technology reshapes the organizational structures around it, and those structures in turn reshape the technology's behavior and effects~\cite{tristr81}. Applied to AI deployment, this reciprocity implies that HITL is not merely a technical override mechanism but a social architecture that must be designed with the same rigor as the agent system itself. An HITL mechanism that is technically sound but socially unusable — requiring excessive cognitive load, providing insufficient context, or violating established organizational workflows — will be circumvented, ignored, or silently disabled by the humans it is meant to govern.

The \bxthree{} Framework distinguishes between HITL-as-ratification (the human approves an already-computed answer) and HITL-as-escalation (the human supplies a resolution when the system has exhausted its provisioned capacity). The Bailout Protocol is firmly in the latter category: it triggers a human intervention only when the \bounds{} layer detects a condition that exceeds the node's autonomous resolution capacity, and it packages the escalation signal with sufficient diagnostic context for the human to act meaningfully. The distinction is consequential. HITL-as-ratification requires the human to evaluate an answer they did not produce, placing cognitive burden on the reviewer without corresponding epistemic authority. HITL-as-escalation positions the human as the authoritative resolution point for an identified epistemic gap — a role that aligns with human judgment strengths and respects institutional accountability structures.

% --------------------------------------------------------------
\subsection{Fail-Safe Design in Safety-Critical Systems}
\label{sec:failsafe}

Fail-safe design — ensuring that a system's failure mode preserves human safety rather than introducing new hazards — is well-established in aerospace, nuclear energy, and industrial process control. Classical fail-safe methodology proceeds by identifying hazardous system states, defining the conditions under which each is entered, and ensuring that the system transitions to a safe state whenever it cannot guarantee continued safe operation. The defining property of a fail-safe system is that its safe state is not contingent on the continued correct functioning of the system itself: the fail-safe mechanism must operate even when the system is in a degraded or erroneous state~\cite{dijkstra1982,wiener1948}.

Dijkstra's work on system correctness established that the reliability of a system is bounded by the reliability of its error-handling paths, not its correct-operation paths. A system that functions correctly under nominal conditions but has poorly specified or unhandled error modes is, in practice, less safe than a system with more limited nominal capability but robust, well-specified error behavior~\cite{dijkstra1982}. This principle applies forcefully to autonomous AI agents: an agent whose correct-operation behavior is well-specified but whose exception-handling is underspecified represents a net safety liability that grows with operational scope.

In multi-agent systems, fail-safe design is complicated by the absence of a single privileged supervisory process. Unlike a safety-critical control system with a defined supervisor, a distributed AI agent network has no global observer capable of assessing overall system state and invoking a global failover. Wienes's cybernetic framework established that control in complex systems is achieved not through central authority but through the propagation of constraint information across feedback loops~\cite{wiener1948}. The watchdog timer pattern — a localized fail-safe mechanism that monitors an individual process and triggers a recovery action when a condition is violated — maps onto this architecture: each node carries its own constraint monitor, and when the monitor detects a violation, it initiates a recovery action that either restores nominal operation or transitions the node to a safe state.

The Bailout Protocol adopts this architecture but adapts it for the epistemic uncertainty inherent in AI systems. A watchdog timer in an aircraft can verify whether a mechanical process is executing; a constraint monitor in an AI agent must evaluate whether the agent's reasoning remains within its provisioned bounds — a detection task that requires fundamentally different mechanisms. The \bxthree{} \bounds{} layer serves as the constraint monitor; the Bailout Protocol is the structured recovery action triggered when constraints cannot be satisfied autonomously.

% --------------------------------------------------------------
\subsection{The BX3 Framework Three-Layer Model}
\label{sec:bx3-model}

The \bxthree{} Framework defines a three-layer cognitive architecture for autonomous AI agent systems: the \purpose{} layer, the \truth{} layer, and the \bounds{} layer, operating over a shared Forensic Ledger maintained by the \fact{} layer. This architecture provides the structural foundation for the Bailout Protocol's accountability and fail-safe guarantees.

The \purpose{} layer is responsible for mission coherence: it evaluates whether proposed actions remain aligned with the objectives specified by the human principal, flagging any action that would constitute mission drift even if technically within the node's operational scope. It operates at the semantic level — intent alignment — and serves as the first line of containment against goal misalignment.

The \truth{} layer enforces factual grounding of all agent outputs against the Forensic Ledger. It verifies that claims made by agents about their own state, the state of other agents, and the state of the external environment are consistent with ledger entries and are not fabricated or hallucinated. Where the \purpose{} layer operates at the semantic level (intent alignment), the \truth{} layer operates at the factual level (correspondence to recorded reality). An action may be intent-aligned but factually unsupported — the \truth{} layer catches this case.

The \bounds{} layer defines and enforces the operational envelope of each node: the set of permitted actions, the conditions under which they may be taken, and the exception propagation rules for conditions that exceed the node's provisioned resolution range. The \bounds{} layer is where the Bailout Protocol is structurally embedded. When it detects a condition that exceeds the node's autonomous capacity, it triggers the protocol's escalation pathway. The \bounds{} layer is not merely a permission system; it is an active constraint-satisfaction engine that monitors operational state in real time and initiates bailout when constraints cannot be satisfied autonomously.

The \fact{} layer maintains the Forensic Ledger — an append-only, cryptographically sealed record of all state transitions, directive issuances, and exception propagations across the agent network. The ledger is the system of record for all accountability assessments: any audit must trace through the ledger to reconstruct the causal chain between a human directive and an outcome. The \fact{} layer does not participate in the Bailout Protocol's decision logic; it records it. This separation of recording from decision ensures that the ledger is available as an authoritative source of truth even when the agent network is in a degraded state, including during an active bailout.

Together, the three layers define a cognitive architecture in which every agent node has a defined scope (\bounds{}), a defined purpose alignment (\purpose{}), a defined factual grounding (\truth{}), and a permanent, auditable record of its operations (\fact{}). The Bailout Protocol operates at the intersection of these three layers: triggered by \bounds{}, its semantics are defined by \truth{} assessment of diagnostic context, and its resolution is recorded in the Forensic Ledger. The protocol is not an add-on to the \bxthree{} Framework — it is the failover mechanism of the \bounds{} layer itself, the structural guarantee that no condition ever leaves the agent network without a human in the loop.

This architecture contrasts with both flat agent systems, in which all cognitive functions are undifferentiated and failure modes are handled homogeneously, and purely hierarchical supervisory systems, in which a central supervisor monitors all subordinates and handles all exceptions. The \bxthree{} model distributes cognitive function across specialized layers while maintaining a single, architecturally defined failover pathway that is not configurable by individual nodes. The Bailout Protocol is the one invariant across the network: the structural commitment that a human is always reachable for any exception a node cannot resolve.
